%%%%%%%%%%%%%%%%
%%% gen 36 %%%
%%%%%%%%%%%%%%%%

\Chapter{36}

\Verse{1}
{
	\hR{וְאֵלֶּה תֹּלְדוֹת עֵשָׂו הוּא אֱדוֹם׃}
}
{
	\eR{And these are the records of Esau: He is Edom.}
}
{
	Here's Esau's begat list.
	
	He's Red.\footnote{Lit. \paleo{אדם} in OT.}
}

\Verse{2}
{
	\hP{עֵשָׂו לָקַח אֶת נָשָׁיו מִבְּנוֹת כְּנָעַן אֶת עָדָה בַּת אֵילוֹן הַחִתִּי וְאֶת אהֳלִיבָמָה בַּת עֲנָה בַּת צִבְעוֹן הַחִוִּי׃}
}
{
	\eP{
		Esau had taken his wives from the daughters of Canaan:
		Adah, daughter of Elon, the Hittite, and Aholibamah,
		daughter of Anah, daughter of Zibeon, the Hivite,
	}
}
{
%	\footnote{
%		Adah, daughter of Elon … Aholibamah, daughter of Anah … Basemath,
%		daughter of Ishmael. These names are in conflict with the names of
%		Esau’s wives given earlier (26:34; 28:9). Rashi gives midrashic
%		explanations for this confusion, Ibn Ezra suggests that they each had
%		two names, and Ramban says that it is possible that Esau’s first two
%		wives died, he married two others, and he renamed his third wife. The
%		character and variety of their solutions indicate that they did not know
%		what the reason was for this confusion in the text. Critical scholars
%		have not solved it in terms of standard source divisions either. It may
%		be a combination of scribal and source reasons. It remains a classic
%		puzzle to be solved.
%	}


	\begingroup
	\scriptsize
	\setlength{\tabcolsep}{0.35em}
	\Table{|p{0.16\linewidth}|p{0.25\linewidth}|p{0.49\linewidth}|}{
		\hline
		\textbf{Proper name} & \textbf{Hebrew} & \textbf{Possible Meaning} \\
		\hline
		Adah & \heb{עָדָה (\emph{ʿĀdāh})} & Usually “ornament/adornment,” from \heb{עדה} in an adornment-related sense. BDB also notes an Arabic comparison possibly suggesting “morning.” \\
		\hline
		Elon & \heb{אֵילוֹן (\emph{ʾÊlôn})} & “Terebinth/oak” or “oak-grove,” related to \heb{אַיִל/אֵלָה} tree vocabulary. \\
		\hline
		Hittite & \heb{חִתִּי (\emph{Ḥittî})} & No persuasive Hebrew lexical etymology. \\
		\hline
		Aholibamah & \heb{אָהֳלִיבָמָה (\emph{ʾĀholîbāmāh})} & “Tent of the high place”: \heb{אֹהֶל} “tent” + \heb{בָּמָה} “high place/elevation.” \\
		\hline
		Anah & \heb{עֲנָה (\emph{ʿĂnāh})} & Could be compared with \heb{ענה}, “answer/respond,” but the proper-name derivation is not certain. \\
		\hline
		Zibeon & \heb{צִבְעוֹן (\emph{Ṣibʿôn})} & Often connected with \heb{צבע}, “color/dye,” hence “colored/variegated”. \\
		\hline
		Hivite & \heb{חִוִּי (\emph{Ḥiwwî})} & No secure Hebrew etymology; proposals involving \heb{חוה} or animal vocabulary are speculative. \\
		\hline
	}
	\endgroup

	\aB{

	Ok so Esau's wives are:

	\begin{enumerate}
	\item Beauty, daughter of Oak Tree the Turkey.
	\item My Tabernacle is Exalted, daughter of Defile daughter of Hyena the Tent Dweller of Wickedness.
	\end{enumerate}
	}


	\begingroup
	\scriptsize
	\setlength{\tabcolsep}{0.35em}
	\Table{|p{0.16\linewidth}|p{0.25\linewidth}|p{0.49\linewidth}|}{
		\hline
		\textbf{Proper name} & \textbf{Hebrew} & \textbf{Possible Meaning} \\
		\hline
		Adah & \heb{עָדָה} (\emph{ʿĀdāh}) & Usually “ornament/adornment,” from \heb{עדה} in an adornment-related sense. BDB also notes an Arabic comparison possibly suggesting “morning.” \\
		\hline
		Elon & \heb{אֵילוֹן} (\emph{ʾÊlôn}) & “Terebinth/oak” or “oak-grove,” related to \heb{אַיִל/אֵלָה} tree vocabulary. \\
		\hline
		Hittite & \heb{חִתִּי} (\emph{Ḥittî}) & No persuasive Hebrew lexical etymology. \\
		\hline
		Aholibamah & \heb{אָהֳלִיבָמָה} (\emph{ʾĀholîbāmāh}) & “Tent of the high place”: \heb{אֹהֶל} “tent” + \heb{בָּמָה} “high place/elevation.” \\
		\hline
		Anah & \heb{עֲנָה} (\emph{ʿĂnāh}) & Could be compared with \heb{ענה}, “answer/respond,” but the proper-name derivation is not certain. \\
		\hline
		Zibeon & \heb{צִבְעוֹן} (\emph{Ṣibʿôn}) & Often connected with \heb{צבע}, “color/dye,” hence “colored/variegated”. \\
		\hline
		Hivite & \heb{חִוִּי} (\emph{Ḥiwwî}) & No secure Hebrew etymology; proposals involving \heb{חוה} or animal vocabulary are speculative. \\
		\hline
	}
	\endgroup

	\aB{

	Ok so Esau's wives are:

	\begin{enumerate}
	\item Beauty, daughter of Oak Tree the Turkey.
	\item My Tabernacle is Exalted, daughter of Defile daughter of Hyena the Tent Dweller of Wickedness.
	\end{enumerate}
	}
}

\Verse{3}
{
	\hP{וְאֶת בָּשְׂמַת בַּת יִשְׁמָעֵאל אֲחוֹת נְבָיוֹת׃}
}
{
	\eP{and Basemath, daughter of Ishmael, sister of Nebaioth.}
}
{
	\begingroup
	\scriptsize
	\setlength{\tabcolsep}{0.35em}
	\Table{|p{0.16\linewidth}|p{0.25\linewidth}|p{0.49\linewidth}|}{
		\hline
		\textbf{Proper name} & \textbf{Hebrew} & \textbf{Formal Hebrew etymology} \\
		\hline
		Basemath & \heb{בָּשְׂמַת} (\emph{Bāśəmat}) & Related to \heb{בֹּשֶׂם}, “balsam, spice, perfume”; thus “fragrance/perfume.” \\
		\hline
		Ishmael & \heb{יִשְׁמָעֵאל} (\emph{Yišmāʿēl}) & From \heb{שמע} “hear” + \heb{אל} “El/God”: “El hears.” \\
		\hline
		Nebaioth & \heb{נְבָיוֹת} (\emph{Nəḇāyôt}) & No secure Biblical Hebrew etymology. A connection with \heb{נוב} “flourish/bear fruit” is possible but weak. \\
		\hline
	}
	\endgroup

	\aB{and Smells Good, daughter of God Hears, sister of High Looking Prophet.}
}

\Verse{4}
{
	\hP{וַתֵּלֶד עָדָה לְעֵשָׂו אֶת אֱלִיפָז וּבָשְׂמַת יָלְדָה אֶת רְעוּאֵל׃}
}
{
	\eP{
		And Adah gave birth for Esau to Eliphaz, and Basemath gave
		birth to Reuel,
	}
}
{
	\begingroup
	\scriptsize
	\setlength{\tabcolsep}{0.35em}
	\Table{|p{0.16\linewidth}|p{0.25\linewidth}|p{0.49\linewidth}|}{
		\hline
		\textbf{Proper name} & \textbf{Hebrew} & \textbf{Formal Hebrew etymology} \\
		\hline
		Adah & \heb{עָדָה} & “Adornment/ornament,” perhaps with an Arabian alternative. \\
		\hline
		Eliphaz & \heb{אֱלִיפַז} (\emph{ʾĔlîpaz}) & Usually parsed \heb{אֵל} + \heb{פָּז}, approximately “El/God is fine gold.” BDB explicitly puts a question mark on this explanation. \\
		\hline
		Basemath & \heb{בָּשְׂמַת} & “Fragrance/perfume.” \\
		\hline
		Reuel & \heb{רְעוּאֵל} (\emph{Rəʿûʾēl}) & Usually \heb{רע/רעה} + \heb{אל}, “friend/companion of El” or “El is a friend.” \\
		\hline
	}
	\endgroup

	And Beauty gave birth to Go{l,}d.

	And Smells Good gave birth to Godfriend.
}

\Verse{5}
{
	\hP{וְאהֳלִיבָמָה יָלְדָה אֶת [יְעוּשׁ] (יעיש) וְאֶת יַעְלָם וְאֶת קֹרַח אֵלֶּה בְּנֵי עֵשָׂו אֲשֶׁר יֻלְּדוּ לוֹ בְּאֶרֶץ כְּנָעַן׃}
}
{
	\eP{
		and Aholibamah gave birth to Jeush and Jalam and Korah.
		These are Esau’s sons, who were born to him in the land of
		Canaan.
	}
}
{
	Qere Ketiv.


	\begingroup
	\scriptsize
	\setlength{\tabcolsep}{0.35em}
	\Table{|p{0.16\linewidth}|p{0.25\linewidth}|p{0.49\linewidth}|}{
		\hline
		\textbf{Proper name} & \textbf{Hebrew} & \textbf{Formal Hebrew etymology} \\
		\hline
		Aholibamah & \heb{אָהֳלִיבָמָה} & “Tent of the high place.” \\
		\hline
		Jeush & \heb{יְעוּשׁ} (\emph{Yəʿûš}) & Possibly “He will gather” or “Assembler.” \\
		\hline
		Jalam & \heb{יַעְלָם} (\emph{Yaʿlām}) & Sometimes related to \heb{עלם}, “hide,” yielding “hidden,” but BDB's comparative material points toward a wider semantic field involving young adulthood/strength. \\
		\hline
		Korah & \heb{קֹרַח} (\emph{Qōraḥ}) & Perhaps from \heb{קרח}, “ice/frost.” \\
		\hline
	}
	\endgroup

	And My Tabernacle is Exalted gave birth to:
	\begin{enumerate}
	\item He's Coming To Help
	\item Hidden Youngman
	\item Baldness
	\end{enumerate}
}

\Verse{6}
{
	\hP{וַיִּקַּח עֵשָׂו אֶת נָשָׁיו וְאֶת בָּנָיו וְאֶת בְּנֹתָיו וְאֶת כּל נַפְשׁוֹת בֵּיתוֹ וְאֶת מִקְנֵהוּ וְאֶת כּל בְּהֶמְתּוֹ וְאֵת כּל קִנְיָנוֹ אֲשֶׁר רָכַשׁ בְּאֶרֶץ כְּנָעַן וַיֵּלֶךְ אֶל אֶרֶץ מִפְּנֵי יַעֲקֹב אָחִיו׃}
}
{
	\eP{
		And Esau took his wives and his sons and his daughters and
		all the persons of his household and his cattle and all of
		his animals and all of his possessions that he had
		acquired in the land of Canaan, and he went to a land,
		from the presence of Jacob, his brother,
	}
}
{
%	\footnote{
%		a land. The name of the land is not given. It may have been omitted due
%		to a scribal error. Or: some understand the verse to mean that he went
%		to a land that was away from Jacob. The Septuagint reads, “he went out
%		of the land of Canaan,” which makes sense.
%	}

	Esau went to a land.

	Away from Jacob.


	Esau went to a land.

	Away from Jacob.


	Esau went to a land.

	Away from Jacob.
}

\Verse{7}
{
	\hP{כִּי הָיָה רְכוּשָׁם רָב מִשֶּׁבֶת יַחְדָּו וְלֹא יָכְלָה אֶרֶץ מְגוּרֵיהֶם לָשֵׂאת אֹתָם מִפְּנֵי מִקְנֵיהֶם׃}
}
{
	\eP{
		because their property was too great for them to live
		together, and the land of their residences was not able to
		suffice them because of their cattle.
	}
}
{
	The town wasn't big enough for the both of them.
}

\Verse{8}
{
	\hP{וַיֵּשֶׁב עֵשָׂו בְּהַר שֵׂעִיר עֵשָׂו הוּא אֱדוֹם׃}
}
{
	\eP{And Esau lived in Mount Seir. Esau: he is Edom.}
}
{
	\begingroup
	\scriptsize
	\setlength{\tabcolsep}{0.35em}
	\Table{|p{0.16\linewidth}|p{0.25\linewidth}|p{0.49\linewidth}|}{
		\hline
		\textbf{Proper name} & \textbf{Hebrew} & \textbf{Possible Meaning} \\
		\hline
		Esau & \heb{עֵשָׂו} & See v.1. \\
		\hline
		Seir & \heb{שֵׂעִיר} (\emph{Śēʿîr}) & Closely related to שֵׂעִיר/שֵׂעָר, “hairy, shaggy / hair.” \\
		\hline
		Edom & \heb{אֱדוֹם} & “Red.” \\
		\hline
	}
	\endgroup

	\aB{And Esau lived on Mount Hair. Esau: he's Red.}
}

\Verse{9}
{
	\hP{וְאֵלֶּה תֹּלְדוֹת עֵשָׂו אֲבִי אֱדוֹם בְּהַר שֵׂעִיר׃}
}
{
	\eP{
		And these are the records of Esau, father of Edom, in
		Mount Seir.
	}
}
{
	These are the records of Esau, father of Red, on Mount Hair.

	\aB{What?}
}

\Verse{10}
{
	\hP{אֵלֶּה שְׁמוֹת בְּנֵי עֵשָׂו אֱלִיפַז בֶּן עָדָה אֵשֶׁת עֵשָׂו רְעוּאֵל בֶּן בָּשְׂמַת אֵשֶׁת עֵשָׂו׃}
}
{
	\eP{
		These are the names of Esau’s sons: Eliphaz, son of Adah,
		Esau’s wife; Reuel, son of Basemath, Esau’s wife.
	}
}
{
	These are the names of Esau's sons: Go(l+1)d, son of Beauty; and Godfriend, son of Smells Good.

	\begin{itemize}
	\item Go{l,}d son of Beauty, Esau's wife.
	\item Godfriend son of Smells Good, Esau's wife.
	\end{itemize}

	\aB{Didn't we just do this?}
}

\Verse{11}
{
	\hP{וַיִּהְיוּ בְּנֵי אֱלִיפָז תֵּימָן אוֹמָר צְפוֹ וְגַעְתָּם וּקְנַז׃}
}
{
	\eP{
		And Eliphaz’s sons were Teman, Omar, Zepho, and Gatam and
		Kenaz.
	}
}
{
	\begingroup
	\scriptsize
	\setlength{\tabcolsep}{0.35em}
	\Table{|p{0.16\linewidth}|p{0.25\linewidth}|p{0.49\linewidth}|}{
		\hline
		\textbf{Proper name} & \textbf{Hebrew} & \textbf{Possible Meaning} \\
		\hline
		Eliphaz & \heb{אֱלִיפַז} & Perhaps “El is fine gold.” \\
		\hline
		Teman & \heb{תֵּימָן} (\emph{Têmān}) & Related to ימין, “right hand,” hence geographically “south” when facing east. \\
		\hline
		Omar & \heb{אוֹמָר} (\emph{ʾÔmār}) & Usually compared with אמר, “say/speak,” perhaps “speaker.” \\
		\hline
		Zepho & \heb{צְפוֹ} (\emph{Ṣəp̄ô}) & Perhaps from צפה, “look/watch,” hence “watcher/gazing.” \\
		\hline
		Gatam & \heb{גַּעְתָּם} (\emph{Gaʿtām}) & No secure Hebrew derivation. \\
		\hline
		Kenaz & \heb{קְנַז} (\emph{Qənaz}) & No secure Hebrew derivation in BDB/NAS lexicography. \\
		\hline
	}
	\endgroup

	Go(l+1)d's sons were:
	\begin{itemize}
	\item Southpaw
	\item He Says
	\item Watchman
	\item Puny / Exhausted / Burnt Valley, and
	\item Hunter.
	\end{itemize}
}

\Verse{12}
{
	\hP{וְתִמְנַע הָיְתָה פִילֶגֶשׁ לֶאֱלִיפַז בֶּן עֵשָׂו וַתֵּלֶד לֶאֱלִיפַז אֶת עֲמָלֵק אֵלֶּה בְּנֵי עָדָה אֵשֶׁת עֵשָׂו׃}
}
{
	\eP{
		And Timna had been a concubine of Eliphaz, son of Esau,
		and she gave birth to Amalek for Eliphaz. These are the
		sons of Adah, Esau’s wife.
	}
}
{
	\begingroup
	\scriptsize
	\setlength{\tabcolsep}{0.35em}
	\Table{|p{0.16\linewidth}|p{0.25\linewidth}|p{0.49\linewidth}|}{
		\hline
		\textbf{Proper name} & \textbf{Hebrew} & \textbf{Possible Meaning} \\
		\hline
		Timna & \heb{תִּמְנָע} (\emph{Timnāʿ}) & Sometimes related to מנע, “withhold/restrain,” perhaps “restraint.” \\
		\hline
		Eliphaz & \heb{אֱלִיפַז} & Perhaps “El is fine gold.” \\
		\hline
		Esau & \heb{עֵשָׂו} & See v.1. \\
		\hline
		Amalek & \heb{עֲמָלֵק} (\emph{ʿĂmālēq}) & No convincing Hebrew derivation. \\
		\hline
		Adah & \heb{עָדָה} & “Adornment,” possibly another comparative explanation. \\
		\hline
	}
	\endgroup

	And Restraint was a concubine of Gold God, son of Esau.

	And she gave birth to Trouble for Gold God.

	These are the sons of Beauty, Esau's wife.
}

\Verse{13}
{
	\hP{וְאֵלֶּה בְּנֵי רְעוּאֵל נַחַת וָזֶרַח שַׁמָּה וּמִזָּה אֵלֶּה הָיוּ בְּנֵי בָשְׂמַת אֵשֶׁת עֵשָׂו׃}
}
{
	\eP{
		And these are the sons of Reuel: Nahath and Zerah, Shammah
		and Mizzeh. These were the sons of Basemath, Esau’s wife.
	}
}
{
	\begingroup
	\scriptsize
	\setlength{\tabcolsep}{0.35em}
	\Table{|p{0.16\linewidth}|p{0.25\linewidth}|p{0.49\linewidth}|}{
		\hline
		\textbf{Proper name} & \textbf{Hebrew} & \textbf{Possible Meaning} \\
		\hline
		Reuel & \heb{רְעוּאֵל} & “Friend/companion of El.” \\
		\hline
		Nahath & \heb{נַחַת} (\emph{Naḥat}) & Identical with/related to נַחַת, “rest, quiet, descent.” \\
		\hline
		Zerah & \heb{זֶרַח} (\emph{Zeraḥ}) & From זרח, “rise/shine,” especially of the sun: “dawn/rising.” \\
		\hline
		Shammah & \heb{שַׁמָּה} (\emph{Šammāh}) & Could be associated with שׁם “there/name” or שמם-related forms, but no single derivation is secure for this name. \\
		\hline
		Mizzeh & \heb{מִזָּה} (\emph{Mizzāh}) & Possibly from a root meaning strong or firm. \\
		\hline
	}
	\endgroup

	Elfriend's sons were:

	\begin{itemize}
	\item Rest
	\item Dawn
	\item Desolation, and
	\item Faints with Fear
	\end{itemize}

	These were the sons of Smells Good, Esau's wife.
}

\Verse{14}
{
	\hP{וְאֵלֶּה הָיוּ בְּנֵי אהֳלִיבָמָה בַת עֲנָה בַּת צִבְעוֹן אֵשֶׁת עֵשָׂו וַתֵּלֶד לְעֵשָׂו אֶת [יְעוּשׁ] (יעיש) וְאֶת יַעְלָם וְאֶת קֹרַח׃}
}
{
	\eP{
		And these were the sons of Aholibamah, daughter of Anah,
		daughter of Zibeon, Esau’s wife: and she gave birth for
		Esau to Jeush and Jalam and Korah.
	}
}
{
	Qere Ketiv.


	\begingroup
	\scriptsize
	\setlength{\tabcolsep}{0.35em}
	\Table{|p{0.16\linewidth}|p{0.25\linewidth}|p{0.49\linewidth}|}{
		\hline
		\textbf{Proper name} & \textbf{Hebrew} & \textbf{Possible Meaning} \\
		\hline
		Aholibamah & \heb{אָהֳלִיבָמָה} & “Tent of the high place.” \\
		\hline
		Anah & \heb{עֲנָה} & Perhaps related to “answer/respond”; comparative Arabian evidence cautions against certainty. \\
		\hline
		Zibeon & \heb{צִבְעוֹן} & Perhaps “variegated/colored” or “hyena.” \\
		\hline
		Esau & \heb{עֵשָׂו} & See v.1. \\
		\hline
		Jeush & \heb{יְעוּשׁ} & Perhaps “he comes to help.” \\
		\hline
		Jalam & \heb{יַעְלָם} & Hebrew “hidden” possible; comparative evidence favors a broader youth/strength semantic field. \\
		\hline
		Korah & \heb{קֹרַח} & Perhaps “baldness”; “ice/frost” is another lexical association. \\
		\hline
	}
	\endgroup


	\aB{My Tabernacle is Exalted, daughter of Defile daughter of Hyena, gave Esau He's-Coming-To-Help, Hidden Youngman, and Baldness.}
}

\Verse{15}
{
	\hP{אֵלֶּה אַלּוּפֵי בְנֵי עֵשָׂו בְּנֵי אֱלִיפַז בְּכוֹר עֵשָׂו אַלּוּף תֵּימָן אַלּוּף אוֹמָר אַלּוּף צְפוֹ אַלּוּף קְנַז׃}
}
{
	\eP{
		These are the chiefs of the children of Esau: the sons of
		Eliphaz, Esau’s firstborn: chief Teman, chief Omar, chief
		Zepho, chief Kenaz,
	}
}
{
	\begingroup
	\scriptsize
	\setlength{\tabcolsep}{0.35em}
	\Table{|p{0.16\linewidth}|p{0.25\linewidth}|p{0.49\linewidth}|}{
		\hline
		\textbf{Proper name} & \textbf{Hebrew} & \textbf{Possible Meaning} \\
		\hline
		Esau & \heb{עֵשָׂו} & See above. \\
		\hline
		Eliphaz & \heb{אֱלִיפַז} & Perhaps “El is fine gold.” \\
		\hline
		Teman & \heb{תֵּימָן} & “South/right-hand.” \\
		\hline
		Omar & \heb{אוֹמָר} & Perhaps “speaker.” \\
		\hline
		Zepho & \heb{צְפוֹ} & Perhaps “watcher/gazing.” \\
		\hline
		Kenaz & \heb{קְנַז} & No secure Hebrew derivation. \\
		\hline
	}
	\endgroup

	\aB{These are the chiefs: Chief Southpaw, Chief Speaker, Chief Watcher, Chief Who-Knows.}
}

\Verse{16}
{
	\hP{אַלּוּף קֹרַח אַלּוּף גַּעְתָּם אַלּוּף עֲמָלֵק אֵלֶּה אַלּוּפֵי אֱלִיפַז בְּאֶרֶץ אֱדוֹם אֵלֶּה בְּנֵי עָדָה׃}
}
{
	\eP{
		chief Korah, chief Gatam, chief Amalek. These are the
		chiefs of Eliphaz in the land of Edom. These are the sons
		of Adah.
	}
}
{
	\begingroup
	\scriptsize
	\setlength{\tabcolsep}{0.35em}
	\Table{|p{0.16\linewidth}|p{0.25\linewidth}|p{0.49\linewidth}|}{
		\hline
		\textbf{Proper name} & \textbf{Hebrew} & \textbf{Possible Meaning} \\
		\hline
		Korah & \heb{קֹרַח} & Perhaps “baldness.” \\
		\hline
		Gatam & \heb{גַּעְתָּם} & No secure derivation. \\
		\hline
		Amalek & \heb{עֲמָלֵק} & No secure derivation. \\
		\hline
		Eliphaz & \heb{אֱלִיפַז} & Perhaps “El is fine gold.” \\
		\hline
		Edom & \heb{אֱדוֹם} & “Red.” \\
		\hline
		Adah & \heb{עָדָה} & “Adornment.” \\
		\hline
	}
	\endgroup

	\aB{Chief Baldness, Chief Who-Knows, Chief Amalek. These are Go{l,}d's chiefs in Redland.}
}

\Verse{17}
{
	\hP{וְאֵלֶּה בְּנֵי רְעוּאֵל בֶּן עֵשָׂו אַלּוּף נַחַת אַלּוּף זֶרַח אַלּוּף שַׁמָּה אַלּוּף מִזָּה אֵלֶּה אַלּוּפֵי רְעוּאֵל בְּאֶרֶץ אֱדוֹם אֵלֶּה בְּנֵי בָשְׂמַת אֵשֶׁת עֵשָׂו׃}
}
{
	\eP{
		And these are the sons of Reuel, son of Esau: chief
		Nahath, chief Zerah, chief Shammah, chief Mizzah. These
		are the chiefs of Reuel in the land of Edom. These are the
		sons of Basemath, Esau’s wife.
	}
}
{
	\begingroup
	\scriptsize
	\setlength{\tabcolsep}{0.35em}
	\Table{|p{0.16\linewidth}|p{0.25\linewidth}|p{0.49\linewidth}|}{
		\hline
		\textbf{Proper name} & \textbf{Hebrew} & \textbf{Possible Meaning} \\
		\hline
		Reuel & \heb{רְעוּאֵל} & “Friend/companion of El.” \\
		\hline
		Esau & \heb{עֵשָׂו} & See above. \\
		\hline
		Nahath & \heb{נַחַת} & “Rest/quiet,” perhaps “settling/descent.” \\
		\hline
		Zerah & \heb{זֶרַח} & “Rising/dawn.” \\
		\hline
		Shammah & \heb{שַׁמָּה} & Uncertain. \\
		\hline
		Mizzah & \heb{מִזָּה} & No secure derivation. \\
		\hline
		Edom & \heb{אֱדוֹם} & “Red.” \\
		\hline
		Basemath & \heb{בָּשְׂמַת} & “Fragrance.” \\
		\hline
	}
	\endgroup

	\aB{Godfriend's chiefs: Chief Rest, Chief Sunrise, Chief Who-Knows, Chief Who-Knows. These are Smells Good's boys in Redland.}
}

\Verse{18}
{
	\hP{וְאֵלֶּה בְּנֵי אהֳלִיבָמָה אֵשֶׁת עֵשָׂו אַלּוּף יְעוּשׁ אַלּוּף יַעְלָם אַלּוּף קֹרַח אֵלֶּה אַלּוּפֵי אהֳלִיבָמָה בַּת עֲנָה אֵשֶׁת עֵשָׂו׃}
}
{
	\eP{
		And these are the sons of Aholibamah, Esau’s wife: chief
		Jeush, chief Jalam, chief Korah. These were the chiefs of
		Aholibamah, daughter of Anah, Esau’s wife.
	}
}
{
	\begingroup
	\scriptsize
	\setlength{\tabcolsep}{0.35em}
	\Table{|p{0.16\linewidth}|p{0.25\linewidth}|p{0.49\linewidth}|}{
		\hline
		\textbf{Proper name} & \textbf{Hebrew} & \textbf{Possible Meaning} \\
		\hline
		Aholibamah & \heb{אָהֳלִיבָמָה} & “Tent of the high place.” \\
		\hline
		Esau & \heb{עֵשָׂו} & See above. \\
		\hline
		Jeush & \heb{יְעוּשׁ} & Perhaps “he comes to help.” \\
		\hline
		Jalam & \heb{יַעְלָם} & Uncertain; Hebrew “hidden” competes with comparative youth/strength readings. \\
		\hline
		Korah & \heb{קֹרַח} & Perhaps “baldness.” \\
		\hline
		Anah & \heb{עֲנָה} & Possibly “answer”; comparative southern parallels. \\
		\hline
	}
	\endgroup

	\aB{My Tabernacle is Exalted's chiefs: Chief He's-Coming-To-Help, Chief Hidden Youngman, Chief Baldness.}
}

\Verse{19}
{
	\hP{אֵלֶּה בְנֵי עֵשָׂו וְאֵלֶּה אַלּוּפֵיהֶם הוּא אֱדוֹם׃}
}
{
	\eP{
		These are the sons of Esau, and these are their chiefs.
		That is Edom.
	}
}
{
	\begingroup
	\scriptsize
	\setlength{\tabcolsep}{0.35em}
	\Table{|p{0.16\linewidth}|p{0.25\linewidth}|p{0.49\linewidth}|}{
		\hline
		\textbf{Proper name} & \textbf{Hebrew} & \textbf{Possible Meaning} \\
		\hline
		Esau & \heb{עֵשָׂו} & See v.1. \\
		\hline
		Edom & \heb{אֱדוֹם} & “Red.” \\
		\hline
	}
	\endgroup

	\aB{Those are Esau's sons, and those are their chiefs. That's Red.}
}

\Verse{20}
{
	\hP{אֵלֶּה בְנֵי שֵׂעִיר הַחֹרִי יֹשְׁבֵי הָאָרֶץ לוֹטָן וְשׁוֹבָל וְצִבְעוֹן וַעֲנָה׃}
}
{
	\eP{
		These are the sons of Seir the Horite, who live in the
		land: Lotan and Shobal and Zibeon and Anah
	}
}
{
	\begingroup
	\scriptsize
	\setlength{\tabcolsep}{0.35em}
	\Table{|p{0.16\linewidth}|p{0.25\linewidth}|p{0.49\linewidth}|}{
		\hline
		\textbf{Proper name} & \textbf{Hebrew} & \textbf{Possible Meaning} \\
		\hline
		Seir & \heb{שֵׂעִיר} & “Hairy/shaggy.” \\
		\hline
		Horite & \heb{חֹרִי} (\emph{Ḥōrî}) & Older explanations connect with חור, “hole/cave,” hence “cave-dweller,” but this is not secure. \\
		\hline
		Lotan & \heb{לוֹטָן} (\emph{Lôṭān}) & Possibly related to לוט, “cover/wrap.” \\
		\hline
		Shobal & \heb{שׁוֹבָל} (\emph{Šôḇāl}) & Sometimes interpreted “flowing.” \\
		\hline
		Zibeon & \heb{צִבְעוֹן} & “Variegated” or perhaps “hyena.” \\
		\hline
		Anah & \heb{עֲנָה} & Perhaps “answer”; not secure. \\
		\hline
	}
	\endgroup

	\aB{These are the sons of Hairy the Horite: Wrapped-Up, Flowing, Hyena, and Defile.}
}

\Verse{21}
{
	\hP{וְדִשׁוֹן וְאֵצֶר וְדִישָׁן אֵלֶּה אַלּוּפֵי הַחֹרִי בְּנֵי שֵׂעִיר בְּאֶרֶץ אֱדוֹם׃}
}
{
	\eP{
		and Dishon and Ezer and Dishan. These are the chiefs of
		the Horites, the children of Seir, in the land of Edom.
	}
}
{
	\begingroup
	\scriptsize
	\setlength{\tabcolsep}{0.35em}
	\Table{|p{0.16\linewidth}|p{0.25\linewidth}|p{0.49\linewidth}|}{
		\hline
		\textbf{Proper name} & \textbf{Hebrew} & \textbf{Possible Meaning} \\
		\hline
		Dishon & \heb{דִּישׁוֹן} (\emph{Dîšôn}) & Corresponds to a Hebrew animal term usually understood as an antelope/gazelle-type animal. \\
		\hline
		Ezer & \heb{אֵצֶר} (\emph{ʾĒṣer}) & From אצר, “store up,” hence “treasure/store.” \\
		\hline
		Dishan & \heb{דִּישָׁן} (\emph{Dîšān}) & Likely variant/related form of Dishon. \\
		\hline
		Horites & \heb{חֹרִי} & Ethnonym of uncertain derivation. \\
		\hline
		Seir & \heb{שֵׂעִיר} & “Hairy/shaggy.” \\
		\hline
		Edom & \heb{אֱדוֹם} & “Red.” \\
		\hline
	}
	\endgroup

	\aB{And Antelope, Treasure, and Antelope-Again. These are the Horite chiefs in Redland.}
}

\Verse{22}
{
	\hP{וַיִּהְיוּ בְנֵי לוֹטָן חֹרִי וְהֵימָם וַאֲחוֹת לוֹטָן תִּמְנָע׃}
}
{
	\eP{
		And the children of Lotan are Hori and Hemam, and Lotan’s
		sister is Timna.
	}
}
{
	\begingroup
	\scriptsize
	\setlength{\tabcolsep}{0.35em}
	\Table{|p{0.16\linewidth}|p{0.25\linewidth}|p{0.49\linewidth}|}{
		\hline
		\textbf{Proper name} & \textbf{Hebrew} & \textbf{Possible Meaning} \\
		\hline
		Lotan & \heb{לוֹטָן} & Perhaps from לוט, “cover.” \\
		\hline
		Hori & \heb{חֹרִי} (\emph{Ḥōrî}) & Same form as the ethnonym “Horite”; exact etymology uncertain. \\
		\hline
		Hemam & \heb{הֵימָם} (\emph{Hêmām}) & No secure Hebrew derivation; Chronicles has the variant Homam. \\
		\hline
		Timna & \heb{תִּמְנָע} & Perhaps “restraint/withholding.” \\
		\hline
	}
	\endgroup

	\aB{Wrapped-Up's children are Horite and Who-Knows, and Wrapped-Up's sister is Restraint.}
}

\Verse{23}
{
	\hP{וְאֵלֶּה בְּנֵי שׁוֹבָל עַלְוָן וּמָנַחַת וְעֵיבָל שְׁפוֹ וְאוֹנָם׃}
}
{
	\eP{
		And these are the children of Shobal: Alvan and Manahath
		and Ebal, Shepho and Onam.
	}
}
{
	\begingroup
	\scriptsize
	\setlength{\tabcolsep}{0.35em}
	\Table{|p{0.16\linewidth}|p{0.25\linewidth}|p{0.49\linewidth}|}{
		\hline
		\textbf{Proper name} & \textbf{Hebrew} & \textbf{Possible Meaning} \\
		\hline
		Shobal & \heb{שׁוֹבָל} & Perhaps “flowing.” \\
		\hline
		Alvan & \heb{עַלְוָן} (\emph{ʿAlwān}) & Older analysis connects with עלה, perhaps “lofty/high.” \\
		\hline
		Manahath & \heb{מָנַחַת} (\emph{Mānaḥat}) & Probably related to נוח/נחת, “rest, resting place.” \\
		\hline
		Ebal & \heb{עֵיבָל} (\emph{ʿÊḇāl}) & No secure derivation. \\
		\hline
		Shepho & \heb{שְׁפוֹ} (\emph{Šəp̄ô}) & Older lexicography gives approximately “bareness/bare.” \\
		\hline
		Onam & \heb{אוֹנָם} (\emph{ʾÔnām}) & Probably related to אוֹן, “strength/vigor,” hence “vigorous.” \\
		\hline
	}
	\endgroup

	\aB{Flowing's children: Lofty, Resting Place, Who-Knows, Bare, and Strong.}
}

\Verse{24}
{
	\hP{וְאֵלֶּה בְנֵי צִבְעוֹן וְאַיָּה וַעֲנָה הוּא עֲנָה אֲשֶׁר מָצָא אֶת הַיֵּמִם בַּמִּדְבָּר בִּרְעֹתוֹ אֶת הַחֲמֹרִים לְצִבְעוֹן אָבִיו׃}
}
{
	\eP{
		And these are the children of Zibeon: Ajah and Anah—that
		is Anah who found the water in the wilderness when he
		tended the asses of Zibeon, his father.
	}
}
{
	\begingroup
	\scriptsize
	\setlength{\tabcolsep}{0.35em}
	\Table{|p{0.16\linewidth}|p{0.25\linewidth}|p{0.49\linewidth}|}{
		\hline
		\textbf{Proper name} & \textbf{Hebrew} & \textbf{Possible Meaning} \\
		\hline
		Zibeon & \heb{צִבְעוֹן} & Perhaps “variegated” or “hyena.” \\
		\hline
		Ajah / Aiah & \heb{אַיָּה} (\emph{ʾAyyāh}) & “Falcon/kite,” identical with an attested bird-name. \\
		\hline
		Anah & \heb{עֲנָה} & Perhaps related to “answer/respond”; uncertain. \\
		\hline
	}
	\endgroup

	\aB{Hyena's children: Falcon and Defile—the Defile who found the water while tending his father's asses.}
}

\Verse{25}
{
	\hP{וְאֵלֶּה בְנֵי עֲנָה דִּשֹׁן וְאהֳלִיבָמָה בַּת עֲנָה׃}
}
{
	\eP{
		And these are the children of Anah: Dishon and Aholibamah,
		daughter of Anah.
	}
}
{
	\begingroup
	\scriptsize
	\setlength{\tabcolsep}{0.35em}
	\Table{|p{0.16\linewidth}|p{0.25\linewidth}|p{0.49\linewidth}|}{
		\hline
		\textbf{Proper name} & \textbf{Hebrew} & \textbf{Possible Meaning} \\
		\hline
		Anah & \heb{עֲנָה} & Perhaps “answer/respond”; uncertain. \\
		\hline
		Dishon & \heb{דִּישׁוֹן} & Probably “antelope/gazelle” by lexical association. \\
		\hline
		Aholibamah & \heb{אָהֳלִיבָמָה} & “Tent of the high place.” \\
		\hline
	}
	\endgroup

	\aB{Defile's children: Antelope and My Tabernacle is Exalted.}
}

\Verse{26}
{
	\hP{וְאֵלֶּה בְּנֵי דִישָׁן חֶמְדָּן וְאֶשְׁבָּן וְיִתְרָן וּכְרָן׃}
}
{
	\eP{
		And these are the children of Dishon: Hemdan and Eshban
		and Ithran and Cheran.
	}
}
{
	\begingroup
	\scriptsize
	\setlength{\tabcolsep}{0.35em}
	\Table{|p{0.16\linewidth}|p{0.25\linewidth}|p{0.49\linewidth}|}{
		\hline
		\textbf{Proper name} & \textbf{Hebrew} & \textbf{Possible Meaning} \\
		\hline
		Dishon & \heb{דִּישׁוֹן} & “Antelope/gazelle.” \\
		\hline
		Hemdan & \heb{חֶמְדָּן} (\emph{Ḥemdān}) & From חמד, “desire/delight,” hence “pleasant/desirable.” \\
		\hline
		Eshban & \heb{אֶשְׁבָּן} (\emph{ʾEšbān}) & No secure Hebrew derivation. \\
		\hline
		Ithran & \heb{יִתְרָן} (\emph{Yitrān}) & From יתר, “excess/abundance/excellence.” \\
		\hline
		Cheran & \heb{כְּרָן} (\emph{Kərān}) & No secure Hebrew derivation. \\
		\hline
	}
	\endgroup

	\aB{Antelope's children: Desirable, Who-Knows, Abundant, and Who-Knows.}
}

\Verse{27}
{
	\hP{אֵלֶּה בְּנֵי אֵצֶר בִּלְהָן וְזַעֲוָן וַעֲקָן׃}
}
{
	\eP{
		These are the children of Ezer: Bilhan and Zaavan and
		Akan.
	}
}
{
	\begingroup
	\scriptsize
	\setlength{\tabcolsep}{0.35em}
	\Table{|p{0.16\linewidth}|p{0.25\linewidth}|p{0.49\linewidth}|}{
		\hline
		\textbf{Proper name} & \textbf{Hebrew} & \textbf{Possible Meaning} \\
		\hline
		Ezer & \heb{אֵצֶר} & “Treasure/store.” \\
		\hline
		Bilhan & \heb{בִּלְהָן} (\emph{Bilhān}) & BDB regards the derivation as uncertain; older proposals connect it with ideas such as timidity. \\
		\hline
		Zaavan & \heb{זַעֲוָן} (\emph{Zaʿăwān}) & Possibly from זוע, “tremble, be agitated,” hence “trembling.” \\
		\hline
		Akan & \heb{עֲקָן} (\emph{ʿĂqān}) & Perhaps connected with a root for twisting/crooking. \\
		\hline
	}
	\endgroup

	\aB{Treasure's children: Who-Knows, Trembler, and Twisted.}
}

\Verse{28}
{
	\hP{אֵלֶּה בְנֵי דִישָׁן עוּץ וַאֲרָן׃}
}
{
	\eP{These are the children of Dishan: Uz and Aran.}
}
{
	\begingroup
	\scriptsize
	\setlength{\tabcolsep}{0.35em}
	\Table{|p{0.16\linewidth}|p{0.25\linewidth}|p{0.49\linewidth}|}{
		\hline
		\textbf{Proper name} & \textbf{Hebrew} & \textbf{Possible Meaning} \\
		\hline
		Dishan & \heb{דִּישָׁן} & Variant/relative of Dishon, probably linked to the antelope term. \\
		\hline
		Uz & \heb{עוּץ} (\emph{ʿÛṣ}) & Sometimes related to עוץ, “counsel/consult.” \\
		\hline
		Aran & \heb{אֲרָן} (\emph{ʾĂrān}) & No secure Hebrew derivation. \\
		\hline
	}
	\endgroup

	\aB{Antelope-Again's children: Counselor and Who-Knows.}
}

\Verse{29}
{
	\hP{אֵלֶּה אַלּוּפֵי הַחֹרִי אַלּוּף לוֹטָן אַלּוּף שׁוֹבָל אַלּוּף צִבְעוֹן אַלּוּף עֲנָה׃}
}
{
	\eP{
		These are the chiefs of the Horites: chief Lotan, chief
		Shobal, chief Zibeon, chief Anah,
	}
}
{
	\begingroup
	\scriptsize
	\setlength{\tabcolsep}{0.35em}
	\Table{|p{0.16\linewidth}|p{0.25\linewidth}|p{0.49\linewidth}|}{
		\hline
		\textbf{Proper name} & \textbf{Hebrew} & \textbf{Possible Meaning} \\
		\hline
		Horites & \heb{חֹרִי} & Ethnonym; “cave-dweller” is an old but uncertain explanation. \\
		\hline
		Lotan & \heb{לוֹטָן} & Perhaps “covered/wrapped.” \\
		\hline
		Shobal & \heb{שׁוֹבָל} & Perhaps “flowing.” \\
		\hline
		Zibeon & \heb{צִבְעוֹן} & Perhaps “variegated” or “hyena.” \\
		\hline
		Anah & \heb{עֲנָה} & Perhaps “answer.” \\
		\hline
	}
	\endgroup

	\aB{These are the Horite chiefs: Chief Wrapped-Up, Chief Flowing, Chief Hyena, Chief Defile.}
}

\Verse{30}
{
	\hP{אַלּוּף דִּשֹׁן אַלּוּף אֵצֶר אַלּוּף דִּישָׁן אֵלֶּה אַלּוּפֵי הַחֹרִי לְאַלֻּפֵיהֶם בְּאֶרֶץ שֵׂעִיר׃}
}
{
	\eP{
		chief Dishon, chief Ezer, chief Dishan. These are the
		chiefs of the Horites by their chiefdoms in the land of
		Seir.
	}
}
{
	\begingroup
	\scriptsize
	\setlength{\tabcolsep}{0.35em}
	\Table{|p{0.16\linewidth}|p{0.25\linewidth}|p{0.49\linewidth}|}{
		\hline
		\textbf{Proper name} & \textbf{Hebrew} & \textbf{Possible Meaning} \\
		\hline
		Dishon & \heb{דִּישׁוֹן} & Probably “antelope/gazelle.” \\
		\hline
		Ezer & \heb{אֵצֶר} & “Treasure/store.” \\
		\hline
		Dishan & \heb{דִּישָׁן} & Related/variant of Dishon. \\
		\hline
		Horites & \heb{חֹרִי} & Uncertain ethnonym. \\
		\hline
		Seir & \heb{שֵׂעִיר} & “Hairy/shaggy.” \\
		\hline
	}
	\endgroup

	\aB{Chief Antelope, Chief Treasure, Chief Antelope-Again. Those are the chiefs in Hairyland.}
}

\Verse{31}
{
	\hJ{וְאֵלֶּה הַמְּלָכִים אֲשֶׁר מָלְכוּ בְּאֶרֶץ אֱדוֹם לִפְנֵי מְלךְ מֶלֶךְ לִבְנֵי יִשְׂרָאֵל׃}
}
{
	\eJ{
		And these are the kings who ruled in the land of Edom
		before a king ruled the children of Israel.
	}
}
{
	\begingroup
	\scriptsize
	\setlength{\tabcolsep}{0.35em}
	\Table{|p{0.16\linewidth}|p{0.25\linewidth}|p{0.49\linewidth}|}{
		\hline
		\textbf{Proper name} & \textbf{Hebrew} & \textbf{Possible Meaning} \\
		\hline
		Edom & \heb{אֱדוֹם} & “Red.” \\
		\hline
		Israel & \heb{יִשְׂרָאֵל} (\emph{Yiśrāʾēl}) & Genesis 32 explains it through שרה, “strive/contend,” + אל: roughly “you have striven with God”; morphologically the name may also be understood with El as subject, “El strives/rules.” \\
		\hline
	}
	\endgroup

	\aB{These are the kings who ruled Redland before Israel even had a king.}
}

\Verse{32}
{
	\hJ{וַיִּמְלֹךְ בֶּאֱדוֹם בֶּלַע בֶּן בְּעוֹר וְשֵׁם עִירוֹ דִּנְהָבָה׃}
}
{
	\eJ{
		And Bela son of Beor ruled in Edom, and the name of his
		city was Dinhabah.
	}
}
{
	\begingroup
	\scriptsize
	\setlength{\tabcolsep}{0.35em}
	\Table{|p{0.16\linewidth}|p{0.25\linewidth}|p{0.49\linewidth}|}{
		\hline
		\textbf{Proper name} & \textbf{Hebrew} & \textbf{Possible Meaning} \\
		\hline
		Bela & \heb{בֶּלַע} (\emph{Belaʿ}) & Identical in consonants with בלע, “swallow/devour”; thus “swallowing/devourer” is an obvious Hebrew association. \\
		\hline
		Beor & \heb{בְּעוֹר} (\emph{Bəʿôr}) & Traditionally connected with בער, “burn,” yielding “burning” or “torch.” \\
		\hline
		Edom & \heb{אֱדוֹם} & “Red.” \\
		\hline
		Dinhabah & \heb{דִּנְהָבָה} (\emph{Dinhāḇāh}) & No secure Hebrew derivation. \\
		\hline
	}
	\endgroup

	\aB{Devourer son of Burning ruled Redland, and his city was Who-Knows.}
}

\Verse{33}
{
	\hJ{וַיָּמת בָּלַע וַיִּמְלֹךְ תַּחְתָּיו יוֹבָב בֶּן זֶרַח מִבּצְרָה׃}
}
{
	\eJ{
		And Bela died, and Jobab son of Zerah from Bozrah ruled in
		his place.
	}
}
{
	\begingroup
	\scriptsize
	\setlength{\tabcolsep}{0.35em}
	\Table{|p{0.16\linewidth}|p{0.25\linewidth}|p{0.49\linewidth}|}{
		\hline
		\textbf{Proper name} & \textbf{Hebrew} & \textbf{Possible Meaning} \\
		\hline
		Bela & \heb{בֶּלַע} & “Swallow/devour” association. \\
		\hline
		Jobab & \heb{יוֹבָב} (\emph{Yôḇāḇ}) & Traditionally connected with יבב, “cry/howl/wail,” hence approximately “howler.” \\
		\hline
		Zerah & \heb{זֶרַח} & “Rising/dawn.” \\
		\hline
		Bozrah & \heb{בָּצְרָה} (\emph{Boṣrāh}) & From בצר, “fortify/enclose”; “fortification/fortress.” \\
		\hline
	}
	\endgroup

	\aB{Devourer died, and Howler son of Sunrise from Fortress ruled instead.}
}

\Verse{34}
{
	\hJ{וַיָּמת יוֹבָב וַיִּמְלֹךְ תַּחְתָּיו חֻשָׁם מֵאֶרֶץ הַתֵּימָנִי׃}
}
{
	\eJ{
		And Jobab died, and Husham of the land of Temani ruled in
		his place.
	}
}
{
	\begingroup
	\scriptsize
	\setlength{\tabcolsep}{0.35em}
	\Table{|p{0.16\linewidth}|p{0.25\linewidth}|p{0.49\linewidth}|}{
		\hline
		\textbf{Proper name} & \textbf{Hebrew} & \textbf{Possible Meaning} \\
		\hline
		Jobab & \heb{יוֹבָב} & Perhaps “howler/wailer.” \\
		\hline
		Husham & \heb{חֻשָׁם} (\emph{Ḥušām}) & Usually connected with חושׁ, “hurry/hasten,” giving something like “hasty one.” \\
		\hline
		Temani / Teman & \heb{תֵּימָנִי / תֵּימָן} & “Southern / south,” ultimately from “right-hand.” \\
		\hline
	}
	\endgroup

	\aB{Howler died, and Hurry-Up from Southland ruled instead.}
}

\Verse{35}
{
	\hJ{וַיָּמת חֻשָׁם וַיִּמְלֹךְ תַּחְתָּיו הֲדַד בֶּן בְּדַד הַמַּכֶּה אֶת מִדְיָן בִּשְׂדֵה מוֹאָב וְשֵׁם עִירוֹ עֲוִית׃}
}
{
	\eJ{
		And Husham died, and Hadad son of Bedad, who struck Midian
		in the field of Moab, ruled in his place, and the name of
		his city was Avith.
	}
}
{
	\begingroup
	\scriptsize
	\setlength{\tabcolsep}{0.35em}
	\Table{|p{0.16\linewidth}|p{0.25\linewidth}|p{0.49\linewidth}|}{
		\hline
		\textbf{Proper name} & \textbf{Hebrew} & \textbf{Possible Meaning} \\
		\hline
		Husham & \heb{חֻשָׁם} & Perhaps “hasty one.” \\
		\hline
		Hadad & \heb{הֲדַד} (\emph{Hadad}) & \textbf{No good Hebrew etymology should be forced.} It is the name of the well-known West Semitic storm deity Hadad/Haddu, corresponding to Akkadian Adad. \\
		\hline
		Bedad & \heb{בְּדַד} (\emph{Bəḏaḏ}) & Probably associated with בדד, “apart/alone.” \\
		\hline
		Midian & \heb{מִדְיָן} (\emph{Midyān}) & Biblical Hebrew can associate it with strife/contention vocabulary, but historical ethnonym uncertain. \\
		\hline
		Moab & \heb{מוֹאָב} (\emph{Môʾāḇ}) & Genesis gives the folk-etymological wordplay “from father” (\emph{mē-ʾāb}); historical origin remains uncertain. \\
		\hline
		Avith & \heb{עֲוִית} (\emph{ʿĂwît}) & Older analysis connects with עוה, “bend/twist,” perhaps “ruin/twisted”; BDB is cautious. \\
		\hline
	}
	\endgroup

	\aB{Hurry-Up died, and Hadad son of Alone—the guy who beat Midian in Moab—ruled instead. His city was Twisted, maybe.}
}

\Verse{36}
{
	\hJ{וַיָּמת הֲדָד וַיִּמְלֹךְ תַּחְתָּיו שַׂמְלָה מִמַּשְׂרֵקָה׃}
}
{
	\eJ{And Hadad died, and Samlah of Masrekah ruled in his place.}
}
{
	\begingroup
	\scriptsize
	\setlength{\tabcolsep}{0.35em}
	\Table{|p{0.16\linewidth}|p{0.25\linewidth}|p{0.49\linewidth}|}{
		\hline
		\textbf{Proper name} & \textbf{Hebrew} & \textbf{Possible Meaning} \\
		\hline
		Hadad & \heb{הֲדַד} & West Semitic storm-deity name, not securely Hebrew-derived. \\
		\hline
		Samlah & \heb{שַׂמְלָה} (\emph{Śamlāh}) & Form resembles שִׂמְלָה/שַׂמְלָה, “garment/cloak.” \\
		\hline
		Masrekah & \heb{מַשְׂרֵקָה} (\emph{Maśrēqāh}) & Often connected with שׂרק / שׂוֹרֵק, “choice vine,” hence “vineyard/place of vines.” BDB itself is less certain about the underlying root. \\
		\hline
	}
	\endgroup

	\aB{Hadad died, and Cloak from Vineyard ruled instead.}
}

\Verse{37}
{
	\hJ{וַיָּמת שַׂמְלָה וַיִּמְלֹךְ תַּחְתָּיו שָׁאוּל מֵרְחֹבוֹת הַנָּהָר׃}
}
{
	\eJ{
		And Samlah died, and Saul of Rehoboth on the river ruled
		in his place.
	}
}
{
	\begingroup
	\scriptsize
	\setlength{\tabcolsep}{0.35em}
	\Table{|p{0.16\linewidth}|p{0.25\linewidth}|p{0.49\linewidth}|}{
		\hline
		\textbf{Proper name} & \textbf{Hebrew} & \textbf{Possible Meaning} \\
		\hline
		Samlah & \heb{שַׂמְלָה} & Perhaps “garment/cloak.” \\
		\hline
		Saul / Shaul & \heb{שָׁאוּל} (\emph{Šāʾûl}) & Qal passive participle of שאל, “ask/request”: “asked for/requested.” \\
		\hline
		Rehoboth & \heb{רְחֹבוֹת} (\emph{Rəḥōḇôt}) & Plural of רְחֹב, “broad place, plaza, street,” from רחב “be broad.” \\
		\hline
	}
	\endgroup

	\aB{Cloak died, and Asked-For from Broad-Plazas-on-the-River ruled instead.}
}

\Verse{38}
{
	\hJ{וַיָּמת שָׁאוּל וַיִּמְלֹךְ תַּחְתָּיו בַּעַל חָנָן בֶּן עַכְבּוֹר׃}
}
{
	\eJ{
		And Saul died, and Baal-hanan son of Achbor ruled in his
		place.
	}
}
{
	\begingroup
	\scriptsize
	\setlength{\tabcolsep}{0.35em}
	\Table{|p{0.16\linewidth}|p{0.25\linewidth}|p{0.49\linewidth}|}{
		\hline
		\textbf{Proper name} & \textbf{Hebrew} & \textbf{Possible Meaning} \\
		\hline
		Saul / Shaul & \heb{שָׁאוּל} & “Asked/requested.” \\
		\hline
		Baal-hanan & \heb{בַּעַל חָנָן} (\emph{Baʿal Ḥānān}) & “Baal is gracious / Baal has shown favor,” from בעל + חנן. \\
		\hline
		Achbor & \heb{עַכְבּוֹר} (\emph{ʿAḵbôr}) & “Mouse.” \\
		\hline
	}
	\endgroup

	\aB{Asked-For died, and Baal-Is-Gracious son of Mouse ruled instead.}
}

\Verse{39}
{
	\hJ{וַיָּמת בַּעַל חָנָן בֶּן עַכְבּוֹר וַיִּמְלֹךְ תַּחְתָּיו הֲדַר וְשֵׁם עִירוֹ פָּעוּ וְשֵׁם אִשְׁתּוֹ מְהֵיטַבְאֵל בַּת מַטְרֵד בַּת מֵי זָהָב׃}
}
{
	\eJ{
		And Baal-hanan son of Achbor died, and Hadar ruled in his
		place, and the name of his city was Pau; and his wife’s
		name was Mehetabel daughter of Matred daughter of Mezahab.
	}
}
{
	\begingroup
	\scriptsize
	\setlength{\tabcolsep}{0.35em}
	\Table{|p{0.16\linewidth}|p{0.25\linewidth}|p{0.49\linewidth}|}{
		\hline
		\textbf{Proper name} & \textbf{Hebrew} & \textbf{Possible Meaning} \\
		\hline
		Baal-hanan & \heb{בַּעַל חָנָן} & “Baal is gracious.” \\
		\hline
		Achbor & \heb{עַכְבּוֹר} & “Mouse.” \\
		\hline
		Hadar & \heb{הֲדַר} (\emph{Hăḏar}) & If this is the correct text, compare הדר, “splendor, majesty, honor.” \\
		\hline
		Pau & \heb{פָּעוּ} (\emph{Pāʿû}) & No secure Hebrew derivation. Chronicles also preserves a variant form, \emph{Pai}. \\
		\hline
		Mehetabel & \heb{מְהֵיטַבְאֵל} (\emph{Məhêṭaḇʾēl}) & Approximately “El/God does good,” “God benefits,” or “whom El makes good.” \\
		\hline
		Matred & \heb{מַטְרֵד} (\emph{Maṭrēḏ}) & Often connected with טרד, “drive/drive away,” perhaps “driving/one who drives.” \\
		\hline
		Mezahab & \heb{מֵי זָהָב} (\emph{Mê Zāhāḇ}) & Read naturally as “waters of gold”: מֵי “waters of” + זָהָב “gold.” \\
		\hline
	}
	\endgroup

	\aB{Baal-Is-Gracious son of Mouse died, and Splendor ruled instead. His wife was God-Does-Good, daughter of Driver, daughter of Waters-of-Gold.}
}

\Verse{40}
{
	\hJ{וְאֵלֶּה שְׁמוֹת אַלּוּפֵי עֵשָׂו לְמִשְׁפְּחֹתָם לִמְקֹמֹתָם בִּשְׁמֹתָם אַלּוּף תִּמְנָע אַלּוּף עַלְוָה אַלּוּף יְתֵת׃}
}
{
	\eJ{
		And these are the names of the chiefs of Esau by their
		families, by their places, by their names: chief Timnah,
		chief Alvah, chief Jetheth,
	}
}
{
	\begingroup
	\scriptsize
	\setlength{\tabcolsep}{0.35em}
	\Table{|p{0.16\linewidth}|p{0.25\linewidth}|p{0.49\linewidth}|}{
		\hline
		\textbf{Proper name} & \textbf{Hebrew} & \textbf{Possible Meaning} \\
		\hline
		Esau & \heb{עֵשָׂו} & See v.1. \\
		\hline
		Timnah & \heb{תִּמְנָע} & Perhaps “restraint/withholding,” from מנע. \\
		\hline
		Alvah & \heb{עַלְוָה} (\emph{ʿAlwāh}) & Older derivation associates it with עלה/height, perhaps “loftiness.” BDB essentially records it as an Edomite proper name rather than confidently translating it. \\
		\hline
		Jetheth & \heb{יְתֵת} (\emph{Yəṯêṯ}) & Derivation explicitly uncertain in the major lexica; older glosses such as “nail” are not secure. \\
		\hline
	}
	\endgroup

	\aB{These are Esau's chiefs by family and place: Chief Restraint, Chief Loftiness, Chief Who-Knows.}
}

\Verse{41}
{
	\hJ{אַלּוּף אהֳלִיבָמָה אַלּוּף אֵלָה אַלּוּף פִּינֹן׃}
}
{
	\eJ{chief Aholibamah, chief Elah, chief Pinon,}
}
{
	\begingroup
	\scriptsize
	\setlength{\tabcolsep}{0.35em}
	\Table{|p{0.16\linewidth}|p{0.25\linewidth}|p{0.49\linewidth}|}{
		\hline
		\textbf{Proper name} & \textbf{Hebrew} & \textbf{Possible Meaning} \\
		\hline
		Aholibamah & \heb{אָהֳלִיבָמָה} & “Tent of the high place.” \\
		\hline
		Elah & \heb{אֵלָה} (\emph{ʾÊlāh}) & “Terebinth” (a tree-name). \\
		\hline
		Pinon & \heb{פִּינֹן} (\emph{Pînōn}) & BDB gives no secure etymology and compares Punon. \\
		\hline
	}
	\endgroup

	\aB{Chief My Tabernacle is Exalted, Chief Terebinth, Chief Who-Knows.}
}

\Verse{42}
{
	\hJ{אַלּוּף קְנַז אַלּוּף תֵּימָן אַלּוּף מִבְצָר׃}
}
{
	\eJ{chief Kenaz, chief Teman, chief Mibzar,}
}
{
	\begingroup
	\scriptsize
	\setlength{\tabcolsep}{0.35em}
	\Table{|p{0.16\linewidth}|p{0.25\linewidth}|p{0.49\linewidth}|}{
		\hline
		\textbf{Proper name} & \textbf{Hebrew} & \textbf{Possible Meaning} \\
		\hline
		Kenaz & \heb{קְנַז} & No secure Hebrew derivation. \\
		\hline
		Teman & \heb{תֵּימָן} & “South/right-hand.” \\
		\hline
		Mibzar & \heb{מִבְצָר} (\emph{Miḇṣār}) & Identical in form with מִבְצָר, “fortress, fortified place, stronghold,” from בצר. \\
		\hline
	}
	\endgroup

	\aB{Chief Who-Knows, Chief Southpaw, Chief Fortress.}
}

\Verse{43}
{
	\hJ{אַלּוּף מַגְדִּיאֵל אַלּוּף עִירָם אֵלֶּה אַלּוּפֵי אֱדוֹם לְמֹשְׁבֹתָם בְּאֶרֶץ אֲחֻזָּתָם הוּא עֵשָׂו אֲבִי אֱדוֹם׃}
}
{
	\eJ{
		chief Magdiel, chief Iram. These are the chiefs of Edom by
		their homes in the land of their possession. That is Esau,
		father of Edom.
	}
}
{
	\begingroup
	\scriptsize
	\setlength{\tabcolsep}{0.35em}
	\Table{|p{0.16\linewidth}|p{0.25\linewidth}|p{0.49\linewidth}|}{
		\hline
		\textbf{Proper name} & \textbf{Hebrew} & \textbf{Possible Meaning} \\
		\hline
		Magdiel & \heb{מַגְדִּיאֵל} (\emph{Magdîʾēl}) & Usually understood with גדל/מגד + אל, something like “El is greatness,” “excellence of El,” or “El makes great.” \\
		\hline
		Iram & \heb{עִירָם} (\emph{ʿÎrām}) & Sometimes associated with עִיר, “city,” but BDB gives no secure derivation. \\
		\hline
		Edom & \heb{אֱדוֹם} & “Red.” \\
		\hline
		Esau & \heb{עֵשָׂו} & Possible “made/formed/rough” association; historically uncertain. \\
		\hline
	}
	\endgroup

	\aB{Chief God-Makes-Great, Chief Who-Knows. Those are Redland's chiefs. That's Esau, father of Red.}
}
